\chapter{Blood Vessels}

\section*{Patient focus}
This chapter explains the examination or clinical problem in accessible language. Patients should use it to prepare questions and understand the purpose, experience, benefits, and limitations of imaging; local instructions and the official report take priority.


\label{ch:ch43_blood_vessels}

Doctors use special imaging to look at your blood vessels---the arteries and veins that carry blood throughout the body. The main techniques are:

\begin{itemize}
  \item \textbf{CT angiography (CTA)}: A CT scan performed after injecting contrast dye into a vein. It produces detailed 3D images of blood vessels and is fast, making it ideal for emergencies such as suspected aortic dissection or pulmonary embolism.
  \item \textbf{MR angiography (MRA)}: Uses MRI technology, sometimes without contrast, to image blood vessels. It avoids radiation but takes longer.
  \item \textbf{Doppler ultrasound}: Uses sound waves to measure blood flow in arteries and veins. It is painless, radiation-free, and is the first test for problems like deep vein thrombosis (DVT) or carotid artery narrowing.
\end{itemize}

These tests help detect blocked or narrowed arteries (which can cause stroke, heart attack, or leg pain), blood clots, aneurysms (dangerous bulges in arteries), and tears in the vessel wall (dissections). The choice of test depends on which vessel is being examined, how urgent the situation is, and your overall health.

---

\section{Modality Selection}
Ultrasound, CTA, MRA, catheter angiography, and nuclear methods are selected according to vessel, urgency, flow, intervention needs, and patient factors. Catheter angiography is invasive and is generally reserved for a defined diagnostic or therapeutic purpose.

\begin{table}[h]
\centering
\caption{Vascular imaging modalities}
\begin{tabular}{p{3cm} p{3.5cm} p{3.5cm} p{3.5cm}}
\hline
\textbf{Modality} & \textbf{Best for} & \textbf{Advantages} & \textbf{Limitations} \\
\hline
Catheter angiography & Intervention (angioplasty, embolization) & Gold standard for lumen, simultaneous treatment & Invasive, access-site complications, radiation \\
CTA & Aorta, PE, trauma, pre-operative planning & Fast, high spatial resolution, widely available & Radiation, iodinated contrast (renal function, allergy) \\
Duplex ultrasound & Carotid stenosis, DVT, peripheral arterial disease & No radiation, bedside, dynamic flow assessment & Operator-dependent, limited by body habitus \\
MRA (contrast-enhanced) & Aorta, renal arteries, run-off & No radiation, good for young patients & Gadolinium (NSF risk in severe renal failure), cost \\
MRA (non-contrast TOF/PC) & Intracranial vessels, renal arteries & No contrast needed & Longer acquisition, motion artefact \\
\hline
\end{tabular}
\end{table}

\begin{pearlbox}[CTA vs MRA: when to choose]
CTA is often preferred in emergencies because it is fast and widely available. MRA may be useful for selected patients and follow-up. Severe kidney impairment or a prior contrast reaction requires agent- and patient-specific planning; eGFR below 30 is not by itself an absolute contraindication to every contrast agent. Noncontrast techniques such as time-of-flight MRA can answer selected vascular questions.
\end{pearlbox}

---

\section{Contrast Timing in CTA}
Precise timing of contrast bolus is critical for diagnostic-quality vascular imaging:
\begin{itemize}
  \item \textbf{Bolus tracking}: A region of interest is placed over the target vessel; scanning begins when a threshold enhancement (typically 100--150 HU) is reached.
  \item \textbf{Test bolus}: A small bolus is injected to determine the time-to-peak enhancement, which is then used for the main injection.
  \item \textbf{Arterial phase}: Chest CTA $\sim 30$ s; abdominal CTA $\sim 35$ s; run-off CTA requires sequential timing or multiphasic acquisition.
  \item \textbf{Venous contamination}: Excessive venous enhancement obscures arterial detail. Shorten delay or use higher injection rates for cleaner arterial images.
\end{itemize}

---

\section{Aortic Disease}

\subsection{Aortic Aneurysm}
An aneurysm is defined as focal aortic dilation $> 1.5$ times the normal diameter:
\begin{itemize}
  \item \textbf{Abdominal aortic aneurysm (AAA)}: Most common; screening with ultrasound in men aged $> 65$ years who have smoked. Repair threshold: $\geq 5.5$ cm (or $\geq 5.0$ cm in women), or rapid growth ($> 1$ cm/year).
  \item \textbf{Thoracic aortic aneurysm}: Marfan syndrome, bicuspid aortic valve, atherosclerosis, and hypertension are risk factors. Repair threshold varies by aetiology: $\geq 5.5$ cm (general), $\geq 5.0$ cm (Marfan), $\geq 5.0$ cm (bicuspid aortic valve).
\end{itemize}

\subsection{Aortic Dissection}
Aortic dissection is a tear in the intima allowing blood to enter the media, creating a false lumen. It is a surgical emergency.

\textbf{Classification systems}:

\begin{table}[h]
\centering
\caption{Aortic dissection classification}
\begin{tabular}{p{3cm} p{4.5cm} p{4.5cm}}
\hline
\textbf{System} & \textbf{Category} & \textbf{Description} \\
\hline
& Type B & Does NOT involve the ascending aorta (medical management unless complicated) \\
& Type II & Limited to ascending aorta \\
& Type IIIa & Descending aorta, distal to left subclavian, above diaphragm \\
& Type IIIb & Descending aorta, distal to diaphragm \\
DeBakey & Type I & Extends from ascending to descending aorta \\
Stanford & Type A & Involves the ascending aorta (surgical emergency) \\
\hline
\end{tabular}
\end{table}

CTA findings: Intimal flap separating true and false lumens. The false lumen is typically larger, with delayed contrast enhancement. Branch vessel involvement (malperfusion) determines complications.

\begin{pearlbox}[Type A vs Type B]
Stanford Type A dissections involve the ascending aorta and require \textbf{emergency surgery} because of the risk of cardiac tamponade, aortic rupture, and coronary malperfusion. Type B dissections are initially managed medically with strict blood pressure and heart rate control, unless complicated by malperfusion, rupture, or refractory pain.
\end{pearlbox}

---

\section{Pulmonary Embolism (PE)}
CT pulmonary angiography (CTPA) is the gold standard for PE diagnosis. Findings include:

\begin{itemize}
  \item \textbf{Acute PE}: Central (saddle) filling defect straddling a bifurcation; or peripheral wedge-shaped defect within a pulmonary artery. Vessel cutoff sign.
  \item \textbf{Chronic PE}: Eccentric, organised thrombus adherent to the vessel wall; webs and bands; mosaic attenuation; RV dilation.
  \item \textbf{Right heart strain}: RV/LV diameter ratio $> 1$ on axial images; interventricular septal bowing; reflux of contrast into the IVC and hepatic veins.
\end{itemize}

\begin{pearlbox}[Wells score and CTPA]
The Wells score stratifies PE probability. CTPA is indicated for patients with intermediate to high probability. A negative CTPA effectively excludes PE (negative predictive value $> 98\%$). Do not order CTPA in very low probability patients---D-dimer is preferred as a first-line test.
\end{pearlbox}

---

\section{Deep Vein Thrombosis (DVT)}
Compression ultrasound is the first-line investigation for suspected DVT:
\begin{itemize}
  \item \textbf{Diagnostic criterion}: Non-compressibility of the vein (a vein that cannot be compressed with the ultrasound probe contains thrombus)
  \item \textbf{Colour Doppler}: Absent or diminished flow; filling defect
  \item \textbf{Common sites}: Popliteal, femoral, common femoral veins (proximal DVT); calf veins (distal DVT---may require follow-up ultrasound to assess propagation)
\end{itemize}

---

\section{Carotid Stenosis}
Carotid Doppler is the screening tool; CTA or MRA provides anatomical detail for surgical planning.

\begin{table}[h]
\centering
\caption{Carotid stenosis grading (NASCET)}
\begin{tabular}{p{3cm} p{4cm} p{4cm}}
\hline
\textbf{Grade} & \textbf{Stenosis \%} & \textbf{Management} \\
\hline
Mild & $< 50\%$ & Medical therapy (statin, antiplatelet) \\
Moderate & $50$--$69\%$ & Consider CEA based on symptoms and risk factors \\
Occlusion & $100\%$ & Medical management; consider EC-IC bypass in selected cases \\
Severe & $70$--$99\%$ & CEA or CAS indicated if symptomatic \\
\hline
\end{tabular}
\end{table}

---

\section{Peripheral Arterial Disease (PAD)}
CTA or MRA from the aortic bifurcation to the feet evaluates PAD. Key findings:
\begin{itemize}
  \item \textbf{Stenosis severity}: Focal vs diffuse disease; calcification (calcified plaque may overestimate stenosis on CTA)
  \item \textbf{TASC classification}: Guides endovascular vs surgical revascularisation
  \item \textbf{Run-off vessels}: Femoropopliteal and infrapopliteal (tibial) disease determines treatment options (bypass graft target vessels)
  \item \textbf{Collateral vessels}: Indicate chronicity and hemodynamic significance
\end{itemize}

---

\section{Venous Anatomy}
Understanding normal venous anatomy is essential for interpreting vascular imaging:
\begin{itemize}
  \item \textbf{Superior vena cava}: Drains upper body into the right atrium
  \item \textbf{Inferior vena cava}: Drains lower body; receives hepatic, renal, and iliac veins
  \item \textbf{Portal venous system}: Splenic vein + SMV $\rightarrow$ portal vein $\rightarrow$ liver. Portal hypertension is diagnosed when the portal vein diameter $> 13$ mm and/or varices are present.
  \item \textbf{Cervical veins}: Internal jugular, external jugular, subclavian---relevant for central line placement and thrombosis
\end{itemize}

---

\section{Interpretation}
Describe patency, stenosis, occlusion, aneurysm, dissection, thrombus, collateralization, runoff, and end-organ effects. Use validated measurement methods and state the reference segment.

\section{Safety}
Assess contrast, kidney function, access-site risk, anticoagulation, pregnancy, and radiation. Acute limb ischemia, mesenteric ischemia, rupture, or large-vessel occlusion requires urgent communication.

\clearpage
